\section{The exact cost along the released infinite trajectory}
\label{sec:infinitecost}
The finite construction above proves no infinite viscous singularity.
There is a precise calculation for the particular proposed extension
that retains the released infinite fields and changes only the forces
by their Laplacians. We record both that calculation and its scope.

The source trajectory in this section means the specific construction
of Alp\"oge--Buckmaster, including their compact corrections. Its
existence and infinite inviscid conclusions are the source's
Theorems 3.2, 7.3, 8.2, and 1.1, not new conclusions of this note.
We use the following exact properties established in those proofs:
at the growth crossing $t_q^a$, activation is complete and
\begin{equation}
 \Theta_q(t_q^a)=-\frac{\nu^2}{\mu}\lambda_q^{-7/8},
 \quad r_q(t_q^a)\geq r_*>\frac12,
 \quad \left|\frac{\Omega_q(t_q^a)}{\Theta_q(t_q^a)}\right|
 \geq\frac{K_q}{2\sigma_{q-1}}.
 \label{eq:sourcecrossingbounds}
\end{equation}
The first is the definition of the stopping time, source (3.8)
and item (i) after (3.11), p.~10. Here $r_*:=1-C_r\epsilon_{\rm sch}>1/2$ is the improved phase-length bound, with the fixed constant $C_r$ from the source estimate; it is not the separate bootstrap constant called $\zeta_-$ in source (3.14). The second is source (3.18),
p.~13 and the tolerance choice on p.~12. The last is the
tracking estimate in the complete growth proof, pp.~29--30,
between (3.63) and (3.64), with the low-order tolerance decreased
within its permitted fixed range. For clarity the extraction of
a numerical factor does not presume equality with the reference
growing eigenline. With $u=\Omega/\Theta$, $u_*=(b/a)^{1/2}$,
and $z=u/u_*$, direct differentiation gives
\[
 z'=\sqrt{ab}(1-z^2)-z(\log u_*)'.
\]
The source seed is chosen on the growing eigenline, so
$z(t_q^{\rm in})=1$. For $q=1$ the eigenline and the ratio
are exact throughout growth, as proved in the finite-stage
section; the following tracking computation applies to later
stages.
If the source error bound is at most $\delta\leq1/4$, its
Riccati barrier gives $z\geq1-\delta$: at $z=1-\delta$,
$z'\geq\sqrt{ab}\{1-(1-\delta)^2-\delta(1-\delta)\}
=\delta\sqrt{ab}>0$. At $z=1+\delta$ the corresponding
upper derivative is $-\delta\sqrt{ab}<0$. A first exit is
therefore impossible. To extract the coefficient ratio, use the
exact source (3.55)--(3.56): during growth the preceding direction
is vertical, $a_q=\mathcal G_q/(K_qr_q)$ and $b_q=K_qd_q$,
where $d_q=r_q\sin(\phi_q-\phi_{q-1})$ and
$\mathcal G_q=\sum_{i<q}w_iA_i\sin(\phi_q-\phi_i)$ includes
the base term with $A_0$ and $w_0=1$. Hence
$u_*=K_q\sqrt{d_qr_q/\mathcal G_q}$. The three source estimates
are $r_q\geq1-\delta$, $d_q/\sin s_q\geq1-\delta$, and
$\mathcal G_q/(\sigma_{q-1}^2\sin s_q)\leq1+\delta$.
Together they give
$u_*\geq(K_q/\sigma_{q-1})(1-\delta)/\sqrt{1+\delta}$.
Their product is at least $K_q/(2\sigma_{q-1})$, because
$(3/4)^2/\sqrt{5/4}>1/2$. This proves the displayed extraction
from the source's explicit tracking estimates without assigning
a sign to an uncontrolled error.

The source also gives
\begin{equation}
 \sigma_{q-1}\leq\nu\sqrt{C_\sigma}\lambda_{q-1}^{1/16},
 \quad \lambda_q=\lambda_{q-1}^{Q_q},\quad Q_q=Q_*+q,
 \quad Q_*\geq200,
 \label{eq:scalesforcost}
\end{equation}
by (3.7),(3.18). The material radius is
$\ell_q^{\rm phys}=\lambda_{q-1}^{-3}/\mu$ for $q\geq2$.
Each older layer and the base are affine on the full new support
(Lemma 4.3, pp.~37--38), while all corrections to the newest
layer vanish on its entire material plateau (Proposition 6.1,
pp.~54--55). These spatial statements are essential to the following
pointwise computation: a bound on a central affine ball alone
would not provide the points used here.

Choose one fixed $s_\dagger\in(-\pi,\pi)$ with
$F''(s_\dagger)\ne0$. Such a point exists for the released
profile. Otherwise $F''$ would vanish on the full period and by
periodicity everywhere, so $F$ would be an affine periodic
function, hence constant; oddness would make it zero, contradicting
$F'(0)=1$. Fix this point once for every layer. At the growth
crossing put
\begin{equation}
 a_q=\frac{s_\dagger}{K_q}p_q,\qquad
 x_q=M_q(t_q^a)a_q.
 \label{eq:testpoint}
\end{equation}
The activation vectors have $|p_q|=1$ and
$\zeta_q=M_q^{-T}p_q$, so
$K_q\zeta_q\cdot x_q=s_\dagger$ exactly. The point belongs
to the plateau as soon as
\[
 \frac{|s_\dagger|}{K_q}<\frac{\ell_q^{\rm phys}}2,
 \quad\hbox{equivalently}\quad
 2|s_\dagger|<\lambda_{q-1}^{Q_q-3}.
\]
This holds for all sufficiently large $q$ by
\eqref{eq:scalesforcost}. The entire neighborhood of this point
in the plateau has constant envelope, zero correction fields,
and affine older fields. No later layer is yet activated at
$t_q^a$. Consequently at this point the Laplacians of the
\emph{full released fields}, not merely those of an isolated
formal wave, are exactly
\begin{align}
 \Delta\theta_E(x_q,t_q^a)
 &=\Theta_q K_q^2r_q^2F''(s_\dagger),\notag\\
 \Delta u_E(x_q,t_q^a)
 &=\Omega_qK_qJ\zeta_q F''(s_\dagger).
 \label{eq:fulllaplacians}
\end{align}
For the second line use
$u_q=\Omega_q J\zeta_qF(s)/(K_qr_q^2)$ on the plateau;
applying the Laplacian contributes $K_q^2r_q^2$ and cancels
the displayed denominator. This verifies both the vector
direction and every frequency and phase-length factor.

Inserting \eqref{eq:sourcecrossingbounds} gives
\begin{align}
 |\kappa\Delta\theta_E(x_q,t_q^a)|
 &\geq\kappa\nu^2\mu r_*^2
 |F''(s_\dagger)|\lambda_q^{9/8},\label{eq:scalarcostlower}\\
 |\varepsilon\Delta u_E(x_q,t_q^a)|
 &\geq\frac{\varepsilon\nu\mu r_*}{2\sqrt{C_\sigma}}
 |F''(s_\dagger)|\lambda_q^{9/8}
 \lambda_{q-1}^{-1/16}.
 \label{eq:velocitycostlower}
\end{align}
For the second line first use
$|\Omega_q|\geq |\Theta_q|K_q/(2\sigma_{q-1})$ and
$|J\zeta_q|=r_q$, and then \eqref{eq:scalesforcost}.
For fixed positive $\kappa$ the first right-hand side tends
to infinity. For fixed positive $\varepsilon$ the second
does also, since
\[
 \lambda_q^{9/8}\lambda_{q-1}^{-1/16}
 =\lambda_{q-1}^{9Q_q/8-1/16}\longrightarrow\infty.
\]
In particular even the case $\kappa=0$, $\varepsilon>0$
has an explicit diverging vector diffusion cost along these
unchanged fields.

The original forces extend smoothly through the source terminal
time and have fixed compact support, and hence are bounded;
this is Proposition 9.4(b),(c), pp.~72--73, proved using
Lemma 9.3, pp.~71--72.
For the specific exact extension
\[
 f_\theta^{\kappa,\varepsilon}=f_\theta^E-\kappa\Delta\theta_E,
 \qquad f_u^{\kappa,\varepsilon}=f_u^E-\varepsilon\Delta u_E,
 \qquad p^{\kappa,\varepsilon}=p_E,
\]
the reverse triangle inequality and
\eqref{eq:scalarcostlower}--\eqref{eq:velocitycostlower}
therefore prove unbounded forces at the corresponding positive
coefficient, before the terminal time. They cannot have a
continuous, and thus cannot have a smooth, terminal extension.
This proves failure of the specified unchanged-pressure map.
There is a stronger, pressure-independent conclusion for these
same fields. Choose $s_\ddagger$ with $F'''(s_\ddagger)\ne0$.
It exists: otherwise $F''$ would be constant, and periodicity
would successively force $F''=F'=0$, contradicting $F'(0)=1$.
Use the plateau point in \eqref{eq:testpoint} with this phase.
The vorticity there equals a spatially constant older vorticity
plus $\Omega_qF'(s)$, so
\[
 \Delta\omega_E=\Omega_qK_q^2r_q^2F'''(s_\ddagger),\qquad
 |\varepsilon\Delta\omega_E|
 \geq\frac{\varepsilon\nu\mu^2r_*^2}{2\sqrt{C_\sigma}}
 |F'''(s_\ddagger)|\lambda_q^{17/8}
 \lambda_{q-1}^{-1/16}\longrightarrow\infty.
\]
For any smooth pressure $p_{\rm new}$ and any vector force
producing the same $u_E,\theta_E$ in the viscous equation,
subtraction of the inviscid momentum equation and application
of curl give the exact identity
\[
 \operatorname{curl}f_{u,\rm new}
 =\operatorname{curl}f_u^E-\varepsilon\Delta\omega_E,
\]
because mixed derivatives of $p_{\rm new}-p_E$ commute.
The source curl force is bounded. The displayed divergence
therefore excludes a terminal extension with bounded first
spatial derivatives for \emph{every} pressure and force that
retain these exact state fields. In fact the failure is local
at the origin and terminal time. Source Theorem 3.2(iv) gives
$\|M_q\|\leq\sqrt{C_\sigma}\lambda_{q-1}^{1/16}$, so the
test points satisfy
\[
 |x_q|\leq\frac{|s_\ddagger|\sqrt{C_\sigma}}\mu
 \lambda_{q-1}^{1/16-Q_q}\longrightarrow0,
 \qquad t_q^a\longrightarrow T_*.
\]
A continuous derivative on a neighborhood of $(0,T_*)$ is
bounded on a smaller compact neighborhood, contradicting the
computed curl divergence there. If a vector field has bounded
first spatial derivatives then its curl is bounded by their
sum, which proves this implication explicitly. This still
does not exclude a modified coupled trajectory, a different
state construction, or viscous blowup; none of those claims
follows from evaluating the unchanged source state.

The completed result of this note is the actual finite forced
stage proved above and below, including its original datum,
control, background feedback, return, hold, support and every
diffusion term. The released infinite inviscid sequence is
reconstructed as source mathematics. A new infinite viscous
sequence with forces smooth through its terminal time is not
established by these calculations.


